% !TeX root = ../../relatorio-projeto-aplicado.tex
% Substitua todas as instrucoes pelo relato do projeto concluido.
\chapter{Introdução e contexto}
\section{Identificação do projeto}
Informe unidade curricular ou vínculo com o Projeto Aplicado, semestre,
período de execução, equipe, orientação e, quando houver, parceiro ou
demandante. Identifique a autoria deste relatório e a contribuição individual
em trabalhos coletivos. O roteiro dialoga com o Manual ESETI em elaboração
\citeind{manualundf}; sua adoção depende das orientações do curso.

\section{Problema e objetivos}
Contextualize a necessidade, o público ou ambiente atendido, a justificativa,
os objetivos e o escopo efetivamente concluído. Separe os objetivos do projeto
das expectativas futuras. Indique o que será apresentado como evidência de
atendimento de cada objetivo.

\chapter{Fundamentação e requisitos}
Apresente os conceitos, referências e requisitos que orientaram o projeto.
Explique restrições e critérios de aceitação. Registre a origem de cada
requisito e diferencie demandas reais de premissas de uma simulação didática.
Cadastre em \texttt{referencias.bib} as fontes efetivamente citadas.

\chapter{Metodologia e desenvolvimento realizado}
\section{Procedimentos e participação}
Descreva os procedimentos executados, materiais, tecnologias, instrumentos e
decisões. Relacione as contribuições dos integrantes às entregas produzidas,
sem atribuir a uma pessoa todo o trabalho coletivo.

\section{Planejado e realizado}
Compare o planejamento com a execução, explicando alterações de escopo,
atrasos, substituições ou atividades não realizadas e seus efeitos. Quando
utilizar cronograma, apresente o percurso realizado no período. Os registros
de acompanhamento podem subsidiar esta análise, mas não substituem o relato final.

\chapter{Solução e entregas}
Descreva o produto, processo, serviço ou material concluído e seu estágio
real. Relacione entregas, arquivos, versões ou outros identificadores às
evidências correspondentes. Informe como acessá-las, conforme as permissões
aplicáveis. Não apresente uma atividade apenas planejada como produto entregue.

\chapter{Resultados e avaliação}
Relacione objetivos, critérios, procedimentos de avaliação e resultados
obtidos. Apresente evidências e discuta resultados negativos ou inconclusivos.
Distinga verificação técnica, validação com pessoas e impacto observado.
Não invente participantes, entrevistas, testes ou métricas para preencher o roteiro.

\chapter{Discussão e limitações}
Analise as decisões, dificuldades, mudanças de escopo e limites dos resultados.
Explique o aprendizado decorrente do trabalho e a contribuição individual.
Indique condições não avaliadas e questões abertas. Se o relatório for
selecionado para um portfólio, preserve a identificação e a rastreabilidade
das evidências, sem presumir critérios de seleção ainda não definidos pelo curso.

\chapter{Conclusão}
Responda aos objetivos com base no que foi realizado e avaliado. Sintetize a
contribuição, os limites e as recomendações. O encerramento do projeto no
semestre pode deixar questões para continuidade; essas questões não devem
substituir a conclusão do escopo efetivamente desenvolvido.
